
\documentclass{article}
\usepackage{colortbl}
\usepackage{makecell}
\usepackage{multirow}
\usepackage{supertabular}

\begin{document}

\newcounter{utterance}

\centering \large Interaction Transcript for game `hot\_air\_balloon', experiment `air\_balloon\_survival\_en\_reasoning off\_hard', episode 4 with Qwen3{-}32B{-}t0.0.
\vspace{24pt}

{ \footnotesize  \setcounter{utterance}{1}
\setlength{\tabcolsep}{0pt}
\begin{supertabular}{c@{$\;$}|p{.15\linewidth}@{}p{.15\linewidth}p{.15\linewidth}p{.15\linewidth}p{.15\linewidth}p{.15\linewidth}}
    \# & $\;$A & \multicolumn{4}{c}{Game Master} & $\;\:$B\\
    \hline

    \theutterance \stepcounter{utterance}  
    & & \multicolumn{4}{p{0.6\linewidth}}{
        \cellcolor[rgb]{0.9,0.9,0.9}{
            \makecell[{{p{\linewidth}}}]{
                \texttt{\tiny{[P1$\langle$GM]}}
                \texttt{You are participating in a collaborative negotiation game.} \\
\\ 
\texttt{Together with another participant, you must agree on a single set of items that will be kept. Each of you has your own view of how much each item matters to you (importance). You do not know how the other participant values the items. Additionally, you are given the effort each item demands.} \\
\texttt{You may only agree on a set if the total effort of the selected items does not exceed a shared limit:} \\
\\ 
\texttt{LIMIT = 2984} \\
\\ 
\texttt{Here are the individual item effort values:} \\
\\ 
\texttt{Item effort = \{"A21": 350, "B93": 276, "A11": 616, "B33": 737, "C08": 531, "B39": 389, "B96": 24, "C85": 125, "A82": 338, "B31": 356, "B20": 143, "A60": 117, "B09": 257, "C81": 921, "B21": 789\}} \\
\\ 
\texttt{Here is your personal view on the importance of each item:} \\
\\ 
\texttt{Item importance values = \{"A21": 138, "B93": 583, "A11": 868, "B33": 822, "C08": 783, "B39": 65, "B96": 262, "C85": 121, "A82": 508, "B31": 780, "B20": 461, "A60": 484, "B09": 668, "C81": 389, "B21": 808\}} \\
\\ 
\texttt{Goal:} \\
\\ 
\texttt{Your goal is to negotiate a shared set of items that benefits you as much as possible (i.e., maximizes total importance to YOU), while staying within the LIMIT. You are not required to make a PROPOSAL in every message {-} you can simply negotiate as well. All tactics are allowed!} \\
\\ 
\texttt{Interaction Protocol:} \\
\\ 
\texttt{You may only use the following structured formats in a message:} \\
\\ 
\texttt{PROPOSAL: \{'A', 'B', 'C', …\}} \\
\texttt{Propose keeping exactly those items.} \\
\texttt{REFUSE: \{'A', 'B', 'C', …\}} \\
\texttt{Explicitly reject opponent's proposal.} \\
\texttt{ARGUMENT: \{'...'\}} \\
\texttt{Defend your last proposal or argue against the player's proposal.} \\
\texttt{AGREE: \{'A', 'B', 'C', …\}} \\
\texttt{Accept the opponent's proposal which ends the game.} \\
\\ 
\\ 
\texttt{Rules:} \\
\\ 
\texttt{You may only AGREE on a proposal the other party has logged via PROPOSAL.} \\
\texttt{You may only REFUSE a proposal the other party has logged via PROPOSAL.} \\
\texttt{Total effort of any PROPOSAL or AGREE set must be ≤ LIMIT.} \\
\texttt{Do NOT reveal your hidden importance scores.} \\
\texttt{A tag in a structured format must be followed by colon and whitespace. The argument must be a python set containing 0 or more strings.} \\
\texttt{So, it must be of the form TAG: \{...\}} \\
\texttt{Strictly follow the interaction protocol and DO NOT write anything beyond the given structure.} \\
\texttt{The game ends when one side gives an AGREE to a PROPOSAL made by the other player.} \\
\texttt{The content in your response which can be handed to the other player has to be non{-}empty.} \\
\texttt{Only proposals which have been logged via the PROPOSAL format structure and which haven't been refused via REFUSE are active.} \\
\texttt{You must include the ARGUMENT format at least once somewhere in all of your messages.} \\
            }
        }
    }
    & \\ \\

    \theutterance \stepcounter{utterance}  
    & \multicolumn{4}{p{0.6\linewidth}}{
        \cellcolor[rgb]{0.9,0.9,0.9}{
            \makecell[{{p{\linewidth}}}]{
                \texttt{\tiny{[P1$\rangle$GM]}}
                \texttt{<think>} \\
\texttt{Okay, let's see. I need to figure out the best set of items to propose that maximizes my total importance without exceeding the effort limit of 2984. First, I should calculate the total effort of all items to see if they can all fit. Let me add them up.} \\
\\ 
\texttt{Looking at the item efforts: A21 is 350, B93 is 276, A11 is 616, B33 is 737, C08 is 531, B39 is 389, B96 is 24, C85 is 125, A82 is 338, B31 is 356, B20 is 143, A60 is 117, B09 is 257, C81 is 921, B21 is 789. Let me sum these up step by step.} \\
\\ 
\texttt{Starting with 0:} \\
\texttt{+350 = 350} \\
\texttt{+276 = 626} \\
\texttt{+616 = 1242} \\
\texttt{+737 = 1979} \\
\texttt{+531 = 2510} \\
\texttt{+389 = 2899} \\
\texttt{+24 = 2923} \\
\texttt{+125 = 3048. Oh, that's over 2984. So all items together are 3048, which is over the limit. So we need to remove some items to get under 2984. The difference is 3048 {-} 2984 = 64. So we need to remove items with a total effort of at least 64.} \\
\\ 
\texttt{Now, my goal is to maximize my own importance. So I need to find a subset of items where the total effort is ≤2984, and the sum of my importance values is as high as possible. Since I can't know the other person's values, I have to focus on my own priorities.} \\
\\ 
\texttt{Looking at my importance values: A21 is 138, B93 is 583, A11 is 868, B33 is 822, C08 is 783, B39} \\
            }
        }
    }
    & & \\ \\

    \theutterance \stepcounter{utterance}  
    & & \multicolumn{4}{p{0.6\linewidth}}{
        \cellcolor[rgb]{0.9,0.9,0.9}{
            \makecell[{{p{\linewidth}}}]{
                \texttt{\tiny{[P1$\langle$GM]}}
                \texttt{Your response did not contain an argument. You must include the structured format ARGUMENT: \{...\} somewhere in your response. Try again} \\
            }
        }
    }
    & \\ \\

    \theutterance \stepcounter{utterance}  
    & \multicolumn{4}{p{0.6\linewidth}}{
        \cellcolor[rgb]{0.9,0.9,0.9}{
            \makecell[{{p{\linewidth}}}]{
                \texttt{\tiny{[P1$\rangle$GM]}}
                \texttt{<think>} \\
\texttt{Okay, let me try to figure this out. The user wants me to propose a set of items that maximizes my own importance while staying under the effort limit. I need to make sure I include an argument in my response.} \\
\\ 
\texttt{First, I'll check the total effort of all items. Adding them up: 350 (A21) + 276 (B93) + 616 (A11) + 737 (B33) + 531 (C08) + 389 (B39) + 24 (B96) + 125 (C85) + 338 (A82) + 356 (B31) + 143 (B20) + 117 (A60) + 257 (B09) + 921 (C81) + 789 (B21). Let me calculate that step by step.} \\
\\ 
\texttt{Starting with 0:} \\
\texttt{350 + 276 = 626} \\
\texttt{626 + 616 = 1242} \\
\texttt{1242 + 737 = 1979} \\
\texttt{1979 + 531 = 2510} \\
\texttt{2510 + 389 = 2899} \\
\texttt{2899 + 24 = 2923} \\
\texttt{2923 + 125 = 3048} \\
\texttt{3048 + 338 = 3386} \\
\texttt{3386 + 356 = 3742} \\
\texttt{3742 + 143 = 3885} \\
\texttt{3885 + 117 = 3999} \\
\texttt{3999 + 257 = 4256} \\
\texttt{4256 + 921 = 5177} \\
\texttt{5177 + 789 = 5966. Wait, that's way over the limit of 2984. So the total effort is 5966, which is way over. So I need to remove items to get under 2984.} \\
\\ 
\texttt{But maybe I made a mistake in adding.} \\
            }
        }
    }
    & & \\ \\

    \theutterance \stepcounter{utterance}  
    & & \multicolumn{4}{p{0.6\linewidth}}{
        \cellcolor[rgb]{0.9,0.9,0.9}{
            \makecell[{{p{\linewidth}}}]{
                \texttt{\tiny{[P1$\langle$GM]}}
                \texttt{Your response did not contain an argument. You must include the structured format ARGUMENT: \{...\} somewhere in your response. Try again} \\
            }
        }
    }
    & \\ \\

    \theutterance \stepcounter{utterance}  
    & \multicolumn{4}{p{0.6\linewidth}}{
        \cellcolor[rgb]{0.9,0.9,0.9}{
            \makecell[{{p{\linewidth}}}]{
                \texttt{\tiny{[P1$\rangle$GM]}}
                \texttt{<think>} \\
\texttt{Okay, let me try again. The user wants me to include an argument in my response. I need to make sure I structure it properly with the ARGUMENT tag.} \\
\\ 
\texttt{First, I need to find a set of items that maximizes my importance without exceeding the effort limit. Let me start by calculating the total effort of all items. Wait, earlier I thought the total was 3048, but when I added again, I got 5966. That can't be right. Let me check again.} \\
\\ 
\texttt{Looking at the item efforts:} \\
\\ 
\texttt{A21: 350} \\
\texttt{B93: 276} \\
\texttt{A11: 616} \\
\texttt{B33: 737} \\
\texttt{C08: 531} \\
\texttt{B39: 389} \\
\texttt{B96: 24} \\
\texttt{C85: 125} \\
\texttt{A82: 338} \\
\texttt{B31: 356} \\
\texttt{B20: 143} \\
\texttt{A60: 117} \\
\texttt{B09: 257} \\
\texttt{C81: 921} \\
\texttt{B21: 789} \\
\\ 
\texttt{Adding them step by step:} \\
\\ 
\texttt{Start with 0.} \\
\\ 
\texttt{350 (A21) + 276 (B93) = 626} \\
\texttt{626 + 616 (A11) = 1242} \\
\texttt{1242 + 737 (B33) = 1979} \\
\texttt{1979 + 531 (C08) = 2510} \\
\texttt{2510 + 389 (B39) = 2899} \\
\texttt{2899 + 24 (B96) = 2923} \\
\texttt{2923 + 125 (C85) = 3048} \\
\texttt{3048 + 338 (A82) = 3386} \\
\texttt{3386 + 356 (B31) = 3742} \\
\texttt{3742 + 143 (B20) = 3885} \\
\texttt{3885 + 117 (A60) = 3999} \\
\texttt{3999 + 257} \\
            }
        }
    }
    & & \\ \\

    \theutterance \stepcounter{utterance}  
    & & & \multicolumn{2}{p{0.3\linewidth}}{
        \cellcolor[rgb]{0.9,0.9,0.9}{
            \makecell[{{p{\linewidth}}}]{
                \texttt{\tiny{[GM$|$GM]}}
                \texttt{abort} \\
            }
        }
    }
    & & \\ \\

    \theutterance \stepcounter{utterance}  
    & & & \multicolumn{2}{p{0.3\linewidth}}{
        \cellcolor[rgb]{0.9,0.9,0.9}{
            \makecell[{{p{\linewidth}}}]{
                \texttt{\tiny{[GM$|$GM]}}
                \texttt{end game} \\
            }
        }
    }
    & & \\ \\

\end{supertabular}
}

\end{document}
